% META: {"id": "TEXP-GRA-ME-001", "chapitre": "TEXP-GRAPHES", "type_objet": "methode", "capacites_codes": ["C1"], "methodes": ["M1"], "mode_creation": "ex_nihilo", "sources_inspiration": [], "status": "needs_review"}
% BEGIN-VERIFY
% from sympy import Matrix
% # graphe : A-B, A-C, B-C, C-D
% M = Matrix([[0,1,1,0],[1,0,1,0],[1,1,0,1],[0,0,1,0]])
% degres = [sum(M.row(i)) for i in range(4)]
% assert degres == [2, 2, 3, 1]
% assert sum(M) == 2*4  # somme des degres = 2 x nombre d'aretes
% END-VERIFY
\begin{fichemethode}{M1}{Utiliser le vocabulaire des graphes}

\quandUtiliser{%
  Utiliser cette méthode lorsque l'énoncé demande de decrire un graphe avec le
  vocabulaire officiel : sommets, arêtes, adjacence, degré, ordre, chaine,
  connexite, graphe complet.
}

\pasApas{%
  \begin{enumerate}
    \item Identifier l'ORDRE (nombre de sommets) et le nombre d'arêtes ; deux
      sommets relies par une arête sont ADJACENTS.
    \item Calculer le DEGRE de chaque sommet : nombre d'arêtes qui lui sont
      incidentes.
    \item Contrôler avec le lemme des poignees de main : la somme des degrés
      vaut deux fois le nombre d'arêtes.
    \item Chaines : une suite de sommets consecutivement adjacents ; sa
      LONGUEUR est son nombre d'arêtes. Le graphe est CONNEXE si deux sommets
      quelconques sont toujours relies par une chaine.
    \item Graphe COMPLET : toutes les paires de sommets sont adjacentes ; il a
      alors $\binom{n}{2} = \frac{n(n-1)}{2}$ arêtes.
  \end{enumerate}
}

\exempleRedige{%
  Soit $G$ le graphe de sommets $A$, $B$, $C$, $D$ et d'arêtes $AB$, $AC$,
  $BC$, $CD$. Decrire $G$.

  \medskip
  \commentaireMarge{Compter les arêtes incidentes a chaque sommet.}%
  $G$ est d'ordre 4 avec 4 arêtes. Degrés : $d(A) = 2$, $d(B) = 2$,
  $d(C) = 3$, $d(D) = 1$. Contrôle :
  $2 + 2 + 3 + 1 = 8 = 2 \times 4$ arêtes.

  \medskip
  $G$ est connexe ($D$ est relie au reste via $C$, et $A$, $B$, $C$ forment un
  triangle), mais non complet : $A$ et $D$ ne sont pas adjacents (ni $B$ et
  $D$). $A - C - D$ est une chaine de longueur 2 reliant $A$ a $D$.
}

\pieges{%
  \begin{itemize}
    \item \textbf{Piège 1.} Confondre ordre (sommets) et taille (arêtes) du
      graphe.
    \item \textbf{Piège 2.} Confondre longueur d'une chaine (nombre d'ARETES)
      et nombre de sommets traverses.
    \item \textbf{Piège 3.} Declarer un graphe complet des qu'il est connexe :
      complet exige TOUTES les adjacences possibles.
  \end{itemize}
}

\verifier{%
  Vérifier le lemme des poignees de main (somme des degrés paire, égale a
  $2 \times$ arêtes), et justifier la connexite en exhibant une chaine entre
  les sommets « eloignes ».
}

\sEntrainer{\refExos{M1}}

\end{fichemethode}
